% ======================================================================
% Section 5 — Results
% Compatible with \documentclass[pdflatex,sn-mathphys-num]{sn-jnl}
% Numbers: Stage 40 canonical 5-seed means (seeds 42–46), sample std ddof=1
% Sources: gsl_stage41/stage41_summary.csv, stage41_improvements.csv,
%          stage41_paired_tests.csv, stage41_summary.json (sparse controls),
%          Stage 44 §7–8 (capacity + 15-min), Stage 46 terminology
% ======================================================================
\section{Results}\label{sec:results}

This section reports the forecasting results under the protocol of
Sect.~\ref{sec:setup}. Unless stated otherwise, every value is the mean RMSE
over the five training seeds $\{42,\dots,46\}$ with the sample standard
deviation in parentheses; graphs are fixed across seeds. Improvements are
relative RMSE reductions with respect to the named reference. Statistical
qualifications follow Sect.~\ref{sec:stats}: with $n=5$ we do not make
definitive significance claims. Interpretation is deferred to the Discussion.

\subsection{The Physical Graph Is Not a Sound Default}\label{sec:res-physical}

Table~\ref{tab:tgcn-main} reports the T-GCN family on both datasets. The
physical road-network graph (T-GCN) is the worst configuration in its family
in every dataset--horizon cell. On Los-loop at PH1, T-GCN attains
$7.88\pm 0.28$ RMSE, whereas the graph-free control T-GCN-NoSpatial attains
$5.25\pm 0.19$ --- a relative RMSE reduction of $33.3\%$ when the physical
graph is removed. The same direction holds at PH2--PH4 on Los-loop
($29.2\%$, $27.0\%$, and $24.1\%$ relative to T-GCN) and at every horizon on
SZ-Taxi ($24.4\%$, $25.1\%$, $25.2\%$, $25.1\%$). Across all eight
dataset$\times$horizon cells, NoSpatial beats Physical in $5/5$ seeds.

The GCN family (Table~\ref{tab:gcn-main}) shows the same reversal: GCN
(physical) attains $8.14\pm 0.11$ on Los-loop PH1 versus $4.88\pm 0.31$ for
GCN-NoSpatial, and remains worse than GCN-NoSpatial at every horizon on both
datasets.

These results establish the graph-free control, not the physical graph, as
the appropriate reference for evaluating learned structure at these operating
points (structural contrast in Fig.~\ref{fig:graphstruct}).

\subsection{A Single Learned Graph Does Not Close the Gap}\label{sec:res-single}

The contemporaneous learned graphs GSL and cGSL
(Sects.~\ref{sec:gsl_contemp}) improve on the physical graph in most
T-GCN cells, but they do not close the gap to the graph-free control.

On Los-loop, T-GCN-GSL and T-GCN-cGSL obtain $5.86\pm 0.21$ and
$5.82\pm 0.23$ at PH1, versus $5.25\pm 0.19$ for NoSpatial: both are
approximately $11\%$ \emph{worse} than removing the graph, while still
reducing error by roughly $26\%$ relative to the physical baseline. The same
ordering --- Physical $>$ single learned graph $>$ NoSpatial in error ---
holds at PH2--PH4 on Los-loop and at every horizon on SZ-Taxi. In no
dataset$\times$horizon cell does either contemporaneous learned graph beat
the graph-free control.

The GCN family is less uniform. GCN-cGSL ($5.76\pm 0.20$ on Los-loop PH1)
beats the physical GCN but still loses to GCN-NoSpatial ($4.88\pm 0.31$).
GCN-GSL is the notable exception: on Los-loop it is often comparable to, or
worse than, the physical GCN ($7.83\pm 0.11$ at PH1) and does not
consistently improve on that baseline. Symmetrization therefore helps the
non-recurrent GCN backbone (cGSL $<$ GSL in error on Los-loop PH1:
$5.76$ vs $7.83$) while the directed and symmetrized variants are nearly
indistinguishable for T-GCN ($5.86$ vs $5.82$). Mechanism-level discussion
of this asymmetry is deferred to Sect.~6.

In summary, replacing the physical adjacency with a single learned
contemporaneous graph recovers only part of the error reduction obtained by
removing the graph entirely. Motivated by this gap, the remainder of the
results examine lag-specific multi-lag graphs and, critically, how those
graphs are consumed.

\subsection{Sparsity and Capacity Controls}\label{sec:res-controls}

Before interpreting any multi-lag gain, we report the matched-sparsity and
capacity controls of Sect.~\ref{sec:controls}. These controls exist
\textbf{only for Los-loop at PH1} and do not establish conclusions for other
cells.

At a matched budget of $30$ directed edges, the random control RandTop30
obtains $6.10\pm 0.12$ and the top-correlation control CorrTop30 obtains
$5.39\pm 0.10$. Both are worse than the graph-free control on the same cell
($5.25\pm 0.19$). By contrast, the DAGMA multi-lag graphs at the same
$30$-edge budget obtain $4.84\pm 0.11$ under fixed per-timestep consumption
(T-GCN-MultiGSL) and $4.49\pm 0.14$ under gated consumption
(T-GCN-MultiGSL-Mix). Table~\ref{tab:controls} summarizes the comparison.

These matched-budget results show that sparsity alone does not explain the
Los-loop PH1 multi-lag gain: two non-learned $30$-edge graphs are worse than
no graph, whereas DAGMA-placed edges with per-lag consumption are better.
The conclusion is scoped to this single experimental cell.

The capacity control (single seed, seed 42, labeled as such) compares
T-GCN-NoSpatial at hidden dimension $64$ ($12{,}672$ parameters,
RMSE $5.14$) with a hidden-dimension-$74$ variant
($16{,}872$ parameters, RMSE $5.14$) that more closely matches the parameter
count of T-GCN-MultiGSL-Mix ($17{,}091$). Increasing capacity on the
graph-free model changes RMSE by less than $0.02$, whereas Mix reaches
$4.49$ at five seeds on the same cell. The Mix gain is therefore not
explained by its additional $4{,}419$ gate parameters alone.

\subsection{Multi-Lag Graphs: Consumption Matters}\label{sec:res-dissociation}

We next compare configurations that use the \emph{same} multi-lag DAGMA
artifacts (Sect.~\ref{sec:gsl_multilag}) under different consumption
mechanisms (Sect.~\ref{sec:consumption}).

On Los-loop at PH1, T-GCN-MultiGSL (fixed cyclic lag assignment; $0$ extra
parameters) obtains $4.84\pm 0.11$, improving on NoSpatial by $7.8\%$.
T-GCN-MultiGSL-Weighted, which learns only three global scalars, obtains
$4.83\pm 0.11$ --- effectively identical to the fixed assignment. The global
lag mixture does not add value beyond the fixed per-timestep rule.

The striking comparison is with GCN-MultiGSL, which consumes the
\emph{identical} three lag graphs as one static union
$\mathbf{A}_{\cup}=\bigvee_k\mathbf{A}_k$ ($28$ of the $30$ lagged edge
slots appear in the union). GCN-MultiGSL obtains $9.78\pm 0.14$ on Los-loop
PH1: worse than the physical GCN ($8.14\pm 0.11$) and more than twice the
error of T-GCN-MultiGSL on the same graphs ($4.84\pm 0.11$). The direction
persists at PH2--PH4 (GCN-MultiGSL $10.05$--$10.27$ versus T-GCN-MultiGSL
$5.45$--$6.26$).

Because the union holds nearly the same edges as the three lag graphs, the
edge set alone cannot account for the difference between the two outcomes
(Fig.~\ref{fig:dissociation}).
At the same time, GCN and T-GCN also differ in backbone architecture, so
this comparison provides evidence that \emph{how the multi-lag structure is
consumed} is a major determinant of whether it is useful; it is not a fully
controlled isolation of consumption as the sole differing factor. On
SZ-Taxi, where the multi-lag graphs retain only $2$ edges, the corresponding
contrast is small in absolute RMSE.

\subsection{Gated Mixing on Los-Loop: The Positive Result}\label{sec:res-mix}

T-GCN-MultiGSL-Mix, which gates the three static lag graphs per node and per
input timestep, is the strongest configuration on Los-loop at every horizon
(Table~\ref{tab:tgcn-main}). Relative to the graph-free control, Mix reduces
RMSE by $14.5\%$, $11.9\%$, $9.2\%$, and $10.9\%$ at PH1--PH4. Relative to
the physical T-GCN baseline, the corresponding reductions are $43.0\%$,
$37.6\%$, $33.7\%$, and $32.3\%$.

At PH1, Mix attains $4.49\pm 0.14$ versus $5.25\pm 0.19$ for NoSpatial. The
paired per-seed difference is positive in $5/5$ seeds at every horizon
(mean paired difference $0.76$, $0.68$, $0.56$, and $0.72$ RMSE at
PH1--PH4; paired-$t$ $p\le 0.0015$ at PH1--PH2 and $p\le 0.00002$ at PH4,
reported for reference under the $n=5$ caveat of Sect.~\ref{sec:stats}).
Mix also improves on the fixed MultiGSL assignment in $5/5$ seeds at every
horizon, by $0.35$--$0.40$ RMSE on Los-loop (approximately $6$--$7\%$
relative). Weighted remains within $0.01$ RMSE of Fixed at every Los-loop
horizon and is treated as a null ablation.

Figure~\ref{fig:perseed} shows the per-seed distributions
for NoSpatial, MultiGSL, Weighted, and Mix on Los-loop PH1--PH4; the means
match Table~\ref{tab:tgcn-main}.

This is the study's main positive result: on Los-loop, lag-specific learned
graphs improve on the graph-free baseline only when they are consumed in a
temporally aligned manner, and the per-node gate adds a further consistent
gain over the fixed assignment.

\subsection{Boundaries: Temporal Resolution and Dataset Dependence}\label{sec:res-boundaries}

\paragraph{SZ-Taxi.}
On SZ-Taxi, every learned variant remains within seed variability of the
graph-free control. At PH1, T-GCN-MultiGSL-Mix obtains $4.119\pm 0.020$
versus $4.120\pm 0.004$ for NoSpatial (difference $\le 0.015$ RMSE at every
horizon; Mix wins $4/5$, $4/5$, $5/5$, and $4/5$ paired seeds at PH1--PH4).
Relative RMSE reductions versus NoSpatial are at most $0.34\%$. The learned
multi-lag graphs on SZ-Taxi retain only $2$ lagged edges (union $2$), versus
$30$ lagged slots (union $28$) on Los-loop. The Los-loop positive result
therefore does not transfer to SZ-Taxi under the present construction.

\paragraph{Temporal-resolution variant of Los-loop.}
Resampling Los-loop to $15$-minute intervals (Sect.~\ref{sec:variant15})
preserves and strengthens the relative improvement of Mix over the graph-free
baseline at matched wall-clock horizons of $15$--$60$ minutes. At PH1 of the
$15$-minute series ($15$ minutes ahead), NoSpatial obtains
$8.600\pm 0.249$, MultiGSL $7.246\pm 0.298$, and Mix $6.240\pm 0.187$
(a $27.4\%$ relative RMSE reduction versus NoSpatial; Mix wins $5/5$ seeds).
The relative reduction persists at PH2--PH4 of the variant
($21.4\%$, $19.8\%$, $16.1\%$). Absolute RMSE magnitudes are not comparable
across sampling intervals, and PH indices of the variant are not equivalent
to PH indices of the $5$-minute series. Full values are collected in
Table~\ref{tab:los15}.

\paragraph{Summary of the boundary.}
Learned multi-lag structure with lag-aligned consumption helps where the
learned graphs are informative (Los-loop at both sampling intervals). Where
the learned graphs are nearly empty (SZ-Taxi), all learned variants collapse
to the graph-free baseline.

% ======================================================================
% Tables
% ======================================================================

\begin{table}[t]
\centering
\caption{T-GCN family: test RMSE (mean, sample standard deviation in
parentheses) over five seeds. Physical: road-network graph.
NoSpatial: identity adjacency. GSL/cGSL: learned contemporaneous graphs.
MultiGSL family: learned multi-lag graphs under fixed, global-weighted, and
gated consumption (Sect.~\ref{sec:method}).}
\label{tab:tgcn-main}
\small
\begin{tabular}{lcccc}
\toprule
Method & PH1 & PH2 & PH3 & PH4 \\
\midrule
\multicolumn{5}{l}{\textit{Los-loop (5\,min)}} \\
T-GCN (Physical) & 7.88 (0.28) & 8.13 (0.09) & 8.37 (0.17) & 8.66 (0.16) \\
T-GCN-NoSpatial & 5.25 (0.19) & 5.76 (0.08) & 6.11 (0.07) & 6.58 (0.12) \\
T-GCN-GSL & 5.86 (0.21) & 6.30 (0.06) & 6.66 (0.12) & 7.04 (0.09) \\
T-GCN-cGSL & 5.82 (0.23) & 6.35 (0.12) & 6.67 (0.11) & 7.09 (0.12) \\
T-GCN-MultiGSL & 4.84 (0.11) & 5.45 (0.14) & 5.94 (0.03) & 6.26 (0.03) \\
T-GCN-MultiGSL-Weighted & 4.83 (0.11) & 5.44 (0.13) & 5.93 (0.03) & 6.25 (0.03) \\
T-GCN-MultiGSL-Mix & \textbf{4.49 (0.14)} & \textbf{5.08 (0.15)} & \textbf{5.55 (0.16)} & \textbf{5.86 (0.07)} \\
\midrule
\multicolumn{5}{l}{\textit{SZ-Taxi (15\,min)}} \\
T-GCN (Physical) & 5.45 (0.11) & 5.55 (0.08) & 5.60 (0.05) & 5.63 (0.06) \\
T-GCN-NoSpatial & 4.12 (0.00) & 4.16 (0.01) & 4.19 (0.01) & 4.23 (0.01) \\
T-GCN-GSL & 4.28 (0.05) & 4.31 (0.03) & 4.33 (0.01) & 4.37 (0.02) \\
T-GCN-cGSL & 4.30 (0.03) & 4.34 (0.02) & 4.38 (0.03) & 4.41 (0.01) \\
T-GCN-MultiGSL & 4.13 (0.02) & 4.16 (0.00) & 4.20 (0.01) & 4.22 (0.00) \\
T-GCN-MultiGSL-Weighted & 4.13 (0.02) & 4.16 (0.00) & 4.20 (0.01) & 4.22 (0.00) \\
T-GCN-MultiGSL-Mix & 4.12 (0.02) & 4.15 (0.01) & 4.18 (0.00) & 4.22 (0.01) \\
\bottomrule
\end{tabular}
\end{table}

\begin{table}[t]
\centering
\caption{GCN family: test RMSE (mean, sample standard deviation in
parentheses) over five seeds. GCN-MultiGSL consumes the union of the three
multi-lag graphs as one static adjacency (Sect.~\ref{sec:consumption}).}
\label{tab:gcn-main}
\small
\begin{tabular}{lcccc}
\toprule
Method & PH1 & PH2 & PH3 & PH4 \\
\midrule
\multicolumn{5}{l}{\textit{Los-loop (5\,min)}} \\
GCN (Physical) & 8.14 (0.11) & 8.54 (0.28) & 8.76 (0.26) & 8.76 (0.29) \\
GCN-NoSpatial & 4.88 (0.31) & 5.60 (0.39) & 6.02 (0.28) & 6.26 (0.26) \\
GCN-GSL & 7.83 (0.11) & 8.48 (0.21) & 8.45 (0.19) & 8.93 (0.12) \\
GCN-cGSL & 5.76 (0.20) & 6.33 (0.32) & 6.68 (0.25) & 6.88 (0.21) \\
GCN-MultiGSL & 9.78 (0.14) & 10.05 (0.16) & 10.24 (0.13) & 10.27 (0.11) \\
\midrule
\multicolumn{5}{l}{\textit{SZ-Taxi (15\,min)}} \\
GCN (Physical) & 5.96 (0.01) & 5.98 (0.00) & 5.99 (0.00) & 6.00 (0.00) \\
GCN-NoSpatial & 4.11 (0.00) & 4.15 (0.00) & 4.19 (0.00) & 4.22 (0.00) \\
GCN-GSL & 4.88 (0.01) & 4.91 (0.00) & 4.93 (0.01) & 4.96 (0.00) \\
GCN-cGSL & 4.64 (0.00) & 4.67 (0.00) & 4.70 (0.00) & 4.73 (0.00) \\
GCN-MultiGSL & 4.82 (0.01) & 4.85 (0.00) & 4.88 (0.01) & 4.90 (0.00) \\
\bottomrule
\end{tabular}
\end{table}

\begin{table}[t]
\centering
\caption{Matched-sparsity controls on Los-loop at PH1 only
(Sect.~\ref{sec:controls}). All graphs use a $30$-edge budget under the
canonical T-GCN protocol; five seeds except the capacity row, which is
single-seed (seed 42) and labeled as such.}
\label{tab:controls}
\small
\begin{tabular}{lccc}
\toprule
Configuration & RMSE & Edges / extra params & vs NoSpatial \\
\midrule
RandTop30 & 6.10 (0.12) & 30 random & $+16.1\%$ \\
CorrTop30 & 5.39 (0.10) & 30 top-$|$Pearson$|$ & $+2.6\%$ \\
T-GCN-NoSpatial & 5.25 (0.19) & identity & --- \\
T-GCN-MultiGSL & 4.84 (0.11) & 30 DAGMA lag slots & $-7.8\%$ \\
T-GCN-MultiGSL-Mix & 4.49 (0.14) & same graphs, gated & $-14.5\%$ \\
\midrule
NoSpatial $h{=}64$ (seed 42) & 5.14 & $12{,}672$ params & --- \\
NoSpatial $h{=}74$ (seed 42) & 5.14 & $16{,}872$ params & $\approx 0\%$ \\
Mix (seed 42) & 4.46 & $17{,}091$ params & $-13.3\%$ \\
\bottomrule
\end{tabular}
\end{table}

\begin{table}[t]
\centering
\caption{Temporal-resolution variant of Los-loop at $15$-minute sampling
(Sect.~\ref{sec:variant15}). PH $k$ denotes $15k$ minutes ahead and is
\emph{not} comparable to PH $k$ of the $5$-minute series. Five seeds.}
\label{tab:los15}
\small
\begin{tabular}{lcccc}
\toprule
Method & PH1 (15\,min) & PH2 (30\,min) & PH3 (45\,min) & PH4 (60\,min) \\
\midrule
T-GCN-NoSpatial & 8.600 (0.249) & 9.311 (0.197) & 10.048 (0.074) & 10.460 (0.196) \\
T-GCN-MultiGSL & 7.246 (0.298) & 8.571 (0.327) & 9.112 (0.077) & 9.756 (0.210) \\
T-GCN-MultiGSL-Mix & \textbf{6.240 (0.187)} & \textbf{7.322 (0.338)} & \textbf{8.063 (0.203)} & \textbf{8.777 (0.326)} \\
\midrule
Mix vs NoSpatial & $-27.4\%$ & $-21.4\%$ & $-19.8\%$ & $-16.1\%$ \\
Mix wins (of 5 seeds) & 5 & 5 & 5 & 5 \\
\bottomrule
\end{tabular}
\end{table}

\begin{figure}[t]
\centering
\includegraphics[width=0.92\linewidth]{figures/fig1_dissociation.pdf}
\caption{Los-loop PH1: the same multi-lag DAGMA artifacts under static-union
consumption (GCN-MultiGSL) versus per-timestep consumption (T-GCN-MultiGSL),
with physical and gated references. Bars are five-seed mean RMSE; whiskers
are sample standard deviations (seeds 42--46). The union retains 28 of the
30 lagged edge slots. Graphs are static statistical-dependency estimates;
the figure does not encode causal influence.}
\label{fig:dissociation}
\end{figure}

\begin{figure}[t]
\centering
\includegraphics[width=0.98\linewidth]{figures/fig2_perseed.pdf}
\caption{Los-loop per-seed test RMSE for the multi-lag consumption ladder
(seeds 42--46). Diamonds mark five-seed means matching
Table~\ref{tab:tgcn-main}. $n=5$; see Sect.~\ref{sec:stats}.}
\label{fig:perseed}
\end{figure}

\begin{figure}[t]
\centering
\includegraphics[width=0.98\linewidth]{figures/fig3_graphstruct.pdf}
\caption{Los-loop graph structure: (a) physical road-network adjacency
($2626$ off-diagonal edges; mean degree $12.69$); (b) union of the three
multi-lag DAGMA graphs after consumer thresholding ($28$ edges; mean degree
$0.14$); (c) and histogram of node degrees. Learned graphs describe
statistical dependency structure, not causal traffic influence.}
\label{fig:graphstruct}
\end{figure}

